\chapter{Cocaine}
 
\section*{Definition and Overview}
Cocaine (benzoylmethylecgonine) is a powerful, highly addictive central nervous system stimulant drug derived from the leaves of the Erythroxylum coca plant native to South America. In clinical medicine, cocaine has historically been used as a local anaesthetic and vasoconstrictor, particularly in nasal, throat, and ophthalmic surgeries, although its use has been largely superseded by safer synthetic alternatives. Globally and in Australia, cocaine is widely used as a recreational illicit drug. It is a class A/Schedule 8 controlled substance. The pharmacology of cocaine involves inhibiting the reuptake of monoamine neurotransmitters (dopamine, norepinephrine, serotonin), leading to intense euphoria, increased energy, and sympathomimetic physiological effects. Illicit cocaine use carries severe risks of acute cardiovascular toxicity, addiction, psychological dependence, and multi-organ damage.
 
\section*{Epidemiology}
Cocaine is the second most commonly used illicit stimulant in Australia after cannabis. According to national drug strategy household surveys, the lifetime prevalence of cocaine use among Australians aged 14 and over is approximately 10--12\%, with annual use hovering around 2--4\%. Cocaine use is most prevalent among young adults, particularly males aged 20 to 39 years. Cocaine-related hospitalisations and emergency department presentations have risen steadily over the past decade, reflecting increased availability, purity, and polydrug use (commonly combining cocaine with alcohol, which forms the highly toxic metabolite cocaethylene). Cocaine use disorders affect individuals across all socioeconomic groups, with significant public health costs associated with emergency care, addiction treatment, and cardiovascular complications.
 
\section*{Aetiology and Pathophysiology}
Cocaine addiction and toxicity arise from its distinct pharmacological profile:
 
\begin{itemize}
    \item \textbf{Dopamine reuptake inhibition:} cocaine binds to the dopamine transporter (DAT), blocking the reuptake of dopamine in the synaptic cleft of the mesolimbic pathway (the brain's reward centre), causing sustained stimulation and euphoria.
    \item \textbf{Sympathetic nervous system stimulation:} by blocking norepinephrine reuptake, cocaine increases circulating catecholamines, causing tachycardia, severe hypertension, vasoconstriction, and pupillary dilation.
    \item \textbf{Sodium channel blockade:} at higher concentrations, cocaine acts as a local anaesthetic by blocking sodium channels, which can alter cardiac conduction and induce arrhythmias.
    \item \textbf{Cocaethylene formation:} when co-ingested with alcohol, the liver metabolises them into cocaethylene, a compound with a longer half-life and greater cardiotoxicity than cocaine alone.
\end{itemize}
 
\section*{Clinical Features}
The clinical presentation of cocaine use can be divided into acute intoxication, chronic effects, and withdrawal:
 
\begin{itemize}
    \item Extreme euphoria, grandiosity, hyperactivity, rapid speech, and hyper-vigilance
    \item Sympathomimetic features: tachycardia, hypertension, diaphoresis, midriasis (dilated pupils), and tremors
    \item Acute adverse symptoms: chest pain, palpitations, severe anxiety, panic attacks, paranoia, and agitation
    \item In severe toxicity: hyperthermia, seizures, cardiac arrhythmias, and acute coronary syndrome (myocardial infarction)
    \item Chronic effects: nasal septal perforation (from snorting), pulmonary barotrauma ("crack lung"), cognitive impairment, and severe dental decay
    \item Withdrawal features ("crash"): dysphoria, anhedonia, fatigue, hypersomnia, increased appetite, and intense drug cravings
\end{itemize}
 
\section*{Differential Diagnosis}
Cocaine toxicity must be differentiated from other causes of acute sympathomimetic syndromes, cardiovascular distress, or acute psychosis:
 
\begin{itemize}
    \item \textbf{Amphetamine and methamphetamine toxicity:} produces a highly similar sympathomimetic presentation, though with a longer duration of action.
    \item \textbf{Pheochromocytoma:} a catecholamine-secreting tumor presenting with severe episodic hypertension, tachycardia, and headache.
    \item \textbf{Acute coronary syndrome (ACS) due to atherosclerosis:} myocardial infarction presenting with chest pain, differentiated by history, age, and toxicology.
    \item \textbf{Primary acute psychosis or mania:} presenting with agitation, paranoia, or delusions, without the marked physical sympathomimetic signs.
    \item \textbf{Thyrotoxicosis:} thyroid storm presenting with tachycardia, hyperthermia, and agitation.
\end{itemize}
 
\section*{When to Seek Medical Care}
Immediate medical evaluation is required if someone is experiencing acute physical or psychiatric symptoms after cocaine use.
 
\textbf{Call 000 (or local emergency services) for:}
\begin{itemize}
    \item Chest pain, pressure, or tightness, or palpitations
    \item Sudden severe headache, weakness, or difficulty speaking (signs of stroke)
    \item Seizures, convulsions, or loss of consciousness
    \item Extreme agitation, aggressive behavior, severe paranoia, or suicidal thoughts
    \item Hyperthermia (dangerously high body temperature) or difficulty breathing
\end{itemize}
 
\section*{Diagnosis and Investigations}
\begin{itemize}
    \item \textbf{Clinical history and physical examination:} focusing on vital signs (blood pressure, heart rate, temperature) and mental state.
    \item \textbf{Electrocardiogram (ECG):} critical to assess for myocardial ischaemia, arrhythmias, or conduction delays.
    \item \textbf{Serum cardiac markers:} troponin levels to evaluate for acute myocardial injury.
    \item \textbf{Urine toxicology screen:} detects cocaine metabolites (e.g., benzoylecgonine) typically for 2--4 days after use.
    \item \textbf{Blood tests:} full blood count, electrolytes, renal function (to check for rhabdomyolysis), and creatine kinase (CK).
    \item \textbf{Chest X-ray:} to investigate chest pain, ruling out lung collapse (pneumothorax) or pulmonary oedema.
\end{itemize}
 
\section*{Management}
Treatment of acute toxicity is primarily supportive, with a focus on controlling cardiovascular and central nervous system hyperactivity:
 
\begin{itemize}
    \item \textbf{Benzodiazepines:} first-line therapy (e.g., diazepam) to manage agitation, tachycardia, hypertension, and seizures.
    \item \textbf{Cardiovascular control:} severe hypertension is treated with vasodilators like phentolamine or nitroprusside. Beta-blockers are generally avoided in acute cocaine toxicity due to the theoretical risk of unopposed alpha-adrenergic stimulation.
    \item \textbf{Hyperthermia management:} aggressive cooling measures (ice packs, cooling blankets) to prevent multi-organ failure.
    \item \textbf{Long-term rehabilitation:} cognitive behavioural therapy (CBT), contingency management, support groups, and addressing underlying mental health issues (dual diagnosis).
\end{itemize}
 
\section*{Complications}
\begin{itemize}
    \item Acute myocardial infarction, life-threatening arrhythmias, and sudden cardiac death
    \item Haemorrhagic or ischaemic stroke due to acute hypertension
    \item Aortic dissection or rupture
    \item Rhabdomyolysis and subsequent acute kidney injury
    \item Nasal septal perforation and chronic sinusitis
    \item Chronic cocaine use disorder, cognitive decline, and psychiatric disorders (e.g., drug-induced psychosis)
\end{itemize}
 
\section*{Prognosis and Outcomes}
The prognosis for acute cocaine toxicity is generally favorable if complications such as myocardial infarction, stroke, or severe hyperthermia are prevented or treated early. Most patients with uncomplicated intoxication recover fully once the drug is cleared from their system. However, the long-term prognosis for cocaine use disorder is challenging, with high rates of relapse. Successful outcomes depend heavily on the individual's engagement in comprehensive behavioral rehabilitation, psychological support, and lifestyle modification.
 
\section*{Prevention and Self-Care}
\begin{itemize}
    \item Complete avoidance of illicit cocaine use is the only definitive way to prevent harm
    \item Avoid co-ingesting cocaine with alcohol or other substances, which dramatically increases cardiovascular risks
    \item Seek professional support, counseling, or rehabilitation programs early if experiencing difficulty controlling cocaine use
    \item Maintain a supportive network of family, friends, and health professionals to assist in recovery and prevent relapse
    \item Engage in stress-reduction techniques and healthy coping mechanisms to manage drug cravings
\end{itemize}
 
\section*{Sources and Further Reading}
The following reputable organisations provide further information for healthcare students and patients seeking reliable, current advice about cocaine and its effects.
 
\begin{itemize}
    \item Alcohol and Drug Foundation (ADF). \textit{Detailed information on cocaine effects, safety, and support.} https://adf.org.au/
    \item Healthdirect Australia. \textit{Cocaine use, health effects, and treatment options.} https://www.healthdirect.gov.au/
    \item Mayo Clinic. \textit{Substance use disorder: diagnosis, treatment, and recovery.} https://www.mayoclinic.org/
    \item MSD Manual. \textit{Cocaine toxicity, clinical features, and management.} https://www.msdmanuals.com/
    \item World Health Organization. \textit{Global health impact of cocaine and other stimulants.} https://www.who.int/
\end{itemize}
